\section*{Chapter \ref{ch:g11:minerals} --- Minerals}

\begin{solution}{exo:g11:minerals:1}
\omterm{def:g11:minerals:cleavage}{Cleavage} is the tendency of a mineral to break along flat planes, parallel to
directions in which the bonds between its atoms are weaker; fracture is
breakage along irregular surfaces, in a mineral with no plane of weakness.
Quartz has no \omterm{def:g11:minerals:cleavage}{cleavage}: it shows a conchoidal fracture, with smooth curved
surfaces like broken glass. A mica has one \omterm{def:g11:minerals:cleavage}{cleavage}, perfect: it peels into
sheets.
\end{solution}

\begin{solution}{exo:g11:minerals:2}
A trace of another element is enough to change the colour of quartz, which
can be clear, white, black, pink, blue or purple; olivine is almost always
green. More reliable: the \omterm{def:g11:minerals:cleavage}{cleavage} (or fracture), the hardness, the \omterm{def:g11:minerals:lustre}{lustre},
the \omterm{def:g11:minerals:habit}{habit} of free crystals, the density.
\end{solution}

\begin{solution}{exo:g11:minerals:3}
(1) Olivine. (2) Biotite. (3) Orthoclase (a feldspar: two \omterm{def:g11:minerals:cleavage}{cleavages}, hardness
6). (4) Quartz (no \omterm{def:g11:minerals:cleavage}{cleavage}, hardness 7).
\end{solution}

\begin{solution}{exo:g11:minerals:4}
(1) Five atoms: one silicon at the centre and four oxygens at the corners of
a tetrahedron. (2) $+4 + 4 \times (-2) = -4$: the group is \ce{SiO4^4-}.
\end{solution}

\begin{solution}{exo:g11:minerals:5}
Olivine: isolated tetrahedra. Pyroxenes: single chains. Amphiboles: double
chains. Micas: sheets. Feldspars and quartz: frameworks.
\end{solution}

\begin{solution}{exo:g11:minerals:6}
(1) $M(\ce{SiO2}) = 28.085 + 2 \times 15.999 = \qty{60.083}{g/mol}$;
oxygen: $31.998/60.083 \approx \qty{53.3}{\percent}$.
(2) $M(\ce{Mg2SiO4}) = 2 \times 24.305 + 28.085 + 4 \times 15.999 =
\qty{140.691}{g/mol}$; oxygen: $63.996/140.691 \approx \qty{45.5}{\percent}$.
(3) The crust, at \qty{46}{\percent}, is of the same order: it is made almost
entirely of oxygen-rich minerals, the \omterm{def:g11:minerals:silicate}{silicates}.
\end{solution}

\begin{solution}{exo:g11:minerals:7}
The white grain has two \omterm{def:g11:minerals:cleavage}{cleavages} and a hardness between 5.5 and 6.5: a
feldspar (6). The grey grain is harder than 6.5 and shows a curved,
conchoidal fracture with no \omterm{def:g11:minerals:cleavage}{cleavage}: quartz (7).
\end{solution}

\begin{solution}{exo:g11:minerals:8}
(1) $4 - 2/2 = 3$ oxygens per silicon: \ce{SiO3}. (2) \ce{(Si,Al)2O6} holds
two silicon (or aluminium) atoms for six oxygens, three per silicon: it is
\ce{SiO3} twice.
\end{solution}

\begin{solution}{exo:g11:minerals:9}
$\rho = 0.55 \times 2.69 + 0.35 \times 3.4 + 0.10 \times 3.32 = 1.480 + 1.190
+ 0.332 \approx \qty{3.00}{g/cm^3}$. The cube has a volume of
\qty{1000}{cm^3}, so a mass of about \qty{3.00}{kg}.
\end{solution}

\begin{solution}{exo:g11:minerals:10}
(1) $0.75 \times 3.32 + 0.25 \times 3.4 = 2.49 + 0.85 = \qty{3.34}{g/cm^3}$.
(2) $(3.34 - 3.00)/3.00 \approx 0.11$: about \qty{11}{\percent} denser.
(3) The basalt and \omterm{def:g11:ocean-ridges:gabbro}{gabbro} of the \omterm{def:g11:ocean-ridges:oceanic-lithosphere}{oceanic crust} are lighter than the
\omterm{def:g11:ocean-ridges:peridotite}{peridotite} of the mantle: they stay on top of it.
\end{solution}

\begin{solution}{exo:g11:minerals:11}
(1) Plagioclase $168/400 = 0.42$, pyroxene $0.33$, olivine $0.15$, holes
$0.10$. (2) The proposition holds for a rock without pores; here a tenth of
the volume is empty. (3) In the solid part, the fractions are $0.42/0.90
\approx 0.467$, $0.367$ and $0.167$:
$\rho_{\text{solid}} \approx 0.467 \times 2.69 + 0.367 \times 3.4 + 0.167
\times 3.32 \approx \qty{3.06}{g/cm^3}$. With its holes, which add volume
and no mass, the rock has the density $0.42 \times 2.69 + 0.33 \times 3.4 +
0.15 \times 3.32 \approx \qty{2.75}{g/cm^3}$, that is $0.9 \times 3.06$.
\end{solution}

\begin{solution}{exo:g11:minerals:12}
(1) Oxygen $464/16.0 = 29.0$~mol; silicon $283/28.1 \approx 10.1$; aluminium
$84/27.0 \approx 3.1$; iron $52/55.8 \approx 0.93$; calcium $46/40.1 \approx
1.15$; magnesium $28/24.3 \approx 1.15$; sodium $23/23.0 = 1.0$; potassium
$15/39.1 \approx 0.38$. (2) In all about $46.8$~mol of atoms, of which
$29.0/46.8 \approx 0.62$: almost two atoms out of three are oxygen.
(3) Large oxygen ions make up most of the atoms, and even more of the
volume: the crust is, at the scale of atoms, a packing of oxygen ions, with
small silicon and metal ions in the spaces between them.
\end{solution}

\begin{solution}{exo:g11:minerals:13}
(1) The mean number of shared corners is $(2 + 3)/2 = 2.5$, so there are
$4 - 2.5/2 = 2.75$ oxygens per silicon. (2) Four silicons for eleven oxygens:
\ce{Si4O11}. (3) \ce{Si7AlO22} holds eight tetrahedra (seven silicon, one
aluminium) for 22 oxygens, that is $22/8 = 2.75$ per tetrahedron, or
\ce{Si4O11} twice. The two \ce{OH} groups are not part of the tetrahedra.
\end{solution}

\begin{solution}{exo:g11:minerals:14}
$0.20 \times 2.62 + 0.10 \times 2.56 + 0.50 \times 2.65 + 0.10 \times 3.09 +
0.10 \times 3.23 = 0.524 + 0.256 + 1.325 + 0.309 + 0.323 \approx
\qty{2.74}{g/cm^3}$, above \qty{2.7}{g/cm^3}; the granite of the example,
richer in quartz and potassium feldspar, is about \qty{2.59}{g/cm^3}, below
\qty{2.6}{g/cm^3}. The rocks of the west are denser because they hold more
plagioclase and more biotite and hornblende, rich in iron and magnesium; the
high-silica granites hold more quartz and potassium feldspar, the lightest
rock-forming minerals.
\end{solution}

\begin{solution}{exo:g11:minerals:15}
(1) Upper crust: $66.6 \times 0.467 \approx \qty{31.1}{\percent}$ silicon;
whole crust: $60.6 \times 0.467 \approx \qty{28.3}{\percent}$. (2) The upper
crust, richer in silica, should contain more quartz and feldspars, and fewer
minerals rich in iron and magnesium; the deeper crust, the reverse. (3) Quartz
and feldspars are the lightest rock-forming minerals: the upper crust should
be lighter than the deeper crust.
\end{solution}

\begin{solution}{pb:g11:minerals:1}
\textbf{1.} Quartz: no \omterm{def:g11:minerals:cleavage}{cleavage}, conchoidal fracture, glassy \omterm{def:g11:minerals:lustre}{lustre}, harder
than the nail (7).

\textbf{2.} Orthoclase, a potassium feldspar: two \omterm{def:g11:minerals:cleavage}{cleavages}, hardness 6 on
the Mohs scale.

\textbf{3.} Plagioclase (here oligoclase), the feldspar with sodium and
calcium.

\textbf{4.} D is biotite, E muscovite: two micas. Their tetrahedra form
sheets, and the bonds between the sheets are weak: they split along this one
perfect \omterm{def:g11:minerals:cleavage}{cleavage}.

\textbf{5.} Quartz and the two feldspars (A, B, C) are frameworks; biotite
(D) contains iron and magnesium.

\textbf{6.} Quartz $150/500 = 0.30$; orthoclase $0.35$; oligoclase $0.25$;
biotite $0.06$; muscovite $0.04$.

\textbf{7.} When the grains are scattered without order, the share of points
on a mineral estimates the share of the area it covers on the face, and that
equals the share of the volume it fills.

\textbf{8.} $0.35 + 0.25 = 0.60$: \qty{60}{\percent} feldspar.

\textbf{9.} By $1/500 = 0.002$, that is \qty{0.2}{\percent}.

\textbf{10.} Yes: feldspar, quartz and mica.

\textbf{11.} $\rho = \phi_1 \rho_1 + \phi_2 \rho_2 + \dots + \phi_k \rho_k$,
with $\phi_1 + \dots + \phi_k = 1$.

\textbf{12.} Quartz $0.30 \times 2.62 = 0.786$; orthoclase $0.35 \times 2.56 =
0.896$; oligoclase $0.25 \times 2.65 = 0.6625$; biotite $0.06 \times 3.09
\approx 0.185$; muscovite $0.04 \times 2.82 \approx 0.113$ (all in
\unit{g/cm^3}).

\textbf{13.} $0.786 + 0.896 + 0.6625 + 0.185 + 0.113 \approx
\qty{2.64}{g/cm^3}$.

\textbf{14.} Orthoclase, at $0.896$, because it is the most abundant. Biotite
and muscovite, at \qty{3.09}{} and \qty{2.82}{g/cm^3}, are denser than the
other minerals and raise the density of the rock above that of quartz and
feldspar alone.

\textbf{15.} \qty{2640}{kg/m^3}: \qty{1}{m^3} of the granite has a mass of
about \qty{2.64}{t}.

\textbf{16.} $100 \times 2.643 \approx \qty{264.3}{g}$.

\textbf{17.} The block displaces \qty{100}{cm^3} of water, of mass
\qty{100}{g}: the balance reads $264.3 - 100 = \qty{164.3}{g}$.

\textbf{18.} $264.3/(264.3 - 164.3) = 264.3/100 \approx \qty{2.64}{g/cm^3}$.

\textbf{19.} Between the two: lighter than the average of \qty{2.69}{g/cm^3},
above the \qty{2.6}{g/cm^3} of the granites richest in silica and potassium.
More plagioclase and more biotite or hornblende, and less quartz and
potassium feldspar, would make a granite denser.

\textbf{20.} The density of the granite is about
\textbf{\qty{2.64}{g/cm^3}}.
\end{solution}
